\chapter*{阅读指引}

同一个在测试中表现良好的模型，用于提醒、可恢复隔离或直接执行不可逆动作，会产生不同后果。
实际使用中的人群、环境、权限与攻击条件也可能不同于测试。本部分从这些后果提出可信要求，说明证据
能够支持的对象和范围，并为条件失效后的限制使用与补救保留依据。

《可信机器学习基础》先把情境转为可检验要求，再分别检查理论、测试及运行记录，并连接
监测、停用和恢复职责。专题章在共同框架内研究不同对象：可靠性检查预测与不确定性，鲁棒性
检查变化后的表现，公平性检查机会与负担；目标对齐与行为约束研究行为怎样满足实际要求，
可解释性研究理解、质疑与补救的证据，隐私保护研究怎样限制数据使用与发布带来的信息泄露，对抗安全
检查主动操纵能否越过保护边界。

这些对象在实际系统中会相遇。拒答可以减少自动处理错误，却也可能增加某组的等待；解释有助审计，
也可能暴露敏感信息；工具权限限制可以减少损害，也可能影响任务完成。阅读时可以沿同一使用情境对照各章，
分别记录要求、证据及失效条件。一项改进提供的支持只属于对应主张，完整使用判断还需检查其他适用
要求是否满足。
